\documentclass[12pt, a4paper]{article}

\usepackage[margin=1in]{geometry}
\usepackage{times}
\usepackage{parskip}
\usepackage{titlesec}
\usepackage{hyperref}
\usepackage{enumitem}
\usepackage{booktabs}
\usepackage{array}
\usepackage{xcolor}
\usepackage{longtable}
\usepackage{listings}
\usepackage{caption}
\usepackage{multirow}

\definecolor{sectioncolor}{RGB}{30, 80, 150}
\definecolor{codebg}{RGB}{245, 245, 245}
\definecolor{codekw}{RGB}{0, 0, 180}
\definecolor{codecomment}{RGB}{80, 120, 80}
\definecolor{codestring}{RGB}{160, 30, 30}

\titleformat{\section}{\large\bfseries\color{sectioncolor}}{}{0em}{}[\titlerule]
\titleformat{\subsection}{\normalsize\bfseries\color{sectioncolor}}{}{0em}{}
\titleformat{\subsubsection}{\normalsize\bfseries}{}{0em}{}

\lstdefinestyle{sqlstyle}{
    backgroundcolor=\color{codebg},
    basicstyle=\ttfamily\small,
    keywordstyle=\color{codekw}\bfseries,
    commentstyle=\color{codecomment},
    stringstyle=\color{codestring},
    breaklines=true,
    frame=single,
    framerule=0.4pt,
    rulecolor=\color{gray!50},
    showstringspaces=false,
    tabsize=4,
    captionpos=b,
    morekeywords={CREATE, TABLE, PRIMARY, KEY, FOREIGN, REFERENCES, NOT, NULL,
                  DEFAULT, AUTO\_INCREMENT, UNIQUE, CHAR, VARCHAR, INT, BIGINT,
                  BOOLEAN, DATETIME, TEXT, LONGTEXT, CURRENT\_TIMESTAMP,
                  CONSTRAINT, INDEX, ON, DELETE, CASCADE, TRUE}
}

\hypersetup{
    colorlinks=true,
    linkcolor=sectioncolor,
    urlcolor=sectioncolor
}

\newcommand{\pk}{\textbf{\textcolor{sectioncolor}{PK}}}
\newcommand{\fk}{\textbf{\textcolor{red!70!black}{FK}}}
\newcommand{\uk}{\textbf{\textcolor{green!50!black}{UK}}}

\title{
    \textbf{GitDB: Database Structure Documentation} \\[0.4em]
    \large Schema Design, Entity Descriptions, and Relationships \\[0.4em]
    \normalsize DISL Mini Project --- SY-11, Lab Batch F-11 \\
    Pune Institute of Computer Technology (PICT), Pune
}

\author{
    Aditya Malode (23325) \and
    Rohan Natekar (23330) \and
    Aarush Oak (23332) \and
    Prathamesh Kinagi (23340)
}

\date{Academic Year 2025--26}

\begin{document}

\maketitle
\tableofcontents
\newpage

% ─────────────────────────────────────────────
\section{Overview}

GitDB uses a dedicated MySQL database named \texttt{gitdb\_meta} to store all
versioning metadata. This design mirrors Git's \texttt{.git/} directory --- the
versioning layer is fully isolated from the target database being versioned,
ensuring zero pollution of user data.

The \texttt{gitdb\_meta} database consists of \textbf{four entities}:

\begin{center}
\begin{tabular}{>{\ttfamily}l l}
\toprule
\textbf{Entity} & \textbf{Purpose} \\
\midrule
USER       & Stores user accounts and authentication credentials \\
REPOSITORY & Tracks each MySQL database registered for versioning \\
COMMITS    & Stores commit history forming a linear Merkle chain \\
SNAPSHOTS  & Stores full schema and data snapshots per table per commit \\
\bottomrule
\end{tabular}
\end{center}

\medskip
\noindent The \texttt{HEAD} pointer --- which tracks the current commit for each
repository --- is stored directly as the \texttt{current\_hash} column inside
the \texttt{REPOSITORY} table, keeping the design normalised to four entities
rather than five.

% ─────────────────────────────────────────────
\section{Entity Descriptions}

% ── USER ──────────────────────────────────────
\subsection{USER}

The \texttt{USER} entity represents individuals who interact with GitDB. Each
user can own multiple repositories and is recorded as the author of every commit
they make. Authentication is handled via a stored password hash; plaintext
passwords are never persisted.

\begin{center}
\begin{tabular}{>{\ttfamily}l >{\ttfamily}l l l}
\toprule
\textbf{Attribute} & \textbf{Type} & \textbf{Constraint} & \textbf{Description} \\
\midrule
user\_id       & INT           & \pk, AUTO\_INCREMENT  & Surrogate primary key \\
username       & VARCHAR(50)   & \uk, NOT NULL         & Unique login handle \\
email          & VARCHAR(100)  & \uk, NOT NULL         & Unique email address \\
password\_hash & VARCHAR(255)  & NOT NULL              & bcrypt/argon2 hash \\
full\_name     & VARCHAR(100)  & NOT NULL              & Display name for commit attribution \\
is\_active     & BOOLEAN       & DEFAULT TRUE          & Soft-delete flag \\
created\_at    & DATETIME      & DEFAULT CURRENT\_TIMESTAMP & Account creation timestamp \\
\bottomrule
\end{tabular}
\end{center}

\subsubsection*{Design Notes}
\begin{itemize}[leftmargin=1.5em]
    \item \texttt{username} and \texttt{email} both carry unique constraints to
    prevent duplicate registrations via either identifier.
    \item \texttt{is\_active} enables soft deletion --- deactivated users retain
    their commit history without orphaning records.
    \item \texttt{full\_name} is used for commit attribution in \texttt{gitdb log}
    output, analogous to \texttt{git config user.name}.
\end{itemize}

% ── REPOSITORY ────────────────────────────────
\subsection{REPOSITORY}

The \texttt{REPOSITORY} entity records each MySQL database that has been
registered with GitDB via \texttt{gitdb init}. It stores the connection
parameters needed to reach the target database, as well as a pointer
(\texttt{current\_hash}) to the commit that HEAD currently refers to.

\begin{center}
\begin{tabular}{>{\ttfamily}l >{\ttfamily}l l l}
\toprule
\textbf{Attribute} & \textbf{Type} & \textbf{Constraint} & \textbf{Description} \\
\midrule
repo\_id        & INT          & \pk, AUTO\_INCREMENT      & Surrogate primary key \\
user\_id        & INT          & \fk\ $\to$ USER, NOT NULL & Owner of this repository \\
repo\_name      & VARCHAR(100) & NOT NULL                  & Human-readable repository name \\
target\_db\_name & VARCHAR(100)& NOT NULL                  & MySQL database being versioned \\
db\_host        & VARCHAR(100) & NOT NULL                  & MySQL server hostname or IP \\
db\_port        & INT          & NOT NULL                  & MySQL server port (default 3306) \\
db\_user        & VARCHAR(50)  & NOT NULL                  & MySQL login username \\
db\_password\_key & VARCHAR(100) & NULL                    & OS keyring service key; password never stored in DB \\
current\_hash   & CHAR(64)     & \fk\ $\to$ COMMITS, NULL  & HEAD pointer (NULL before first commit) \\
created\_at     & DATETIME     & DEFAULT CURRENT\_TIMESTAMP & Repository registration timestamp \\
\bottomrule
\end{tabular}
\end{center}

\subsubsection*{Design Notes}
\begin{itemize}[leftmargin=1.5em]
    \item \texttt{current\_hash} serves as the HEAD pointer, eliminating the
    need for a separate \texttt{HEAD} table. It is \texttt{NULL} for a freshly
    initialised repository with no commits yet.
    \item The MySQL connection password is \emph{never} stored in \texttt{gitdb\_meta}
    or \texttt{.gitdb/config.json}. Instead, the password is stored in the OS-level
    keyring (via the \texttt{keyring} library) and only the keyring service key
    (e.g.\ \texttt{gitdb:repo\_1}) is persisted in \texttt{db\_password\_key}.
    \item \texttt{current\_hash} is updated on every successful
    \texttt{gitdb commit} and \texttt{gitdb checkout} operation.
\end{itemize}

% ── COMMITS ───────────────────────────────────
\subsection{COMMITS}

The \texttt{COMMITS} entity stores the versioning history of a repository. Each
row represents one commit. Commits are linked via \texttt{parent\_hash} to form
a \textbf{linear Merkle chain} --- since GitDB does not support branching,
each commit has at most one parent and at most one child, making the history
a simple linked list.

\begin{center}
\begin{tabular}{>{\ttfamily}l >{\ttfamily}l l l}
\toprule
\textbf{Attribute} & \textbf{Type} & \textbf{Constraint} & \textbf{Description} \\
\midrule
hash         & CHAR(64)     & \pk                           & SHA-256 content hash \\
repo\_id     & INT          & \fk\ $\to$ REPOSITORY, NOT NULL & Repository this commit belongs to \\
parent\_hash & CHAR(64)     & \fk\ $\to$ COMMITS, \uk, NULL & Hash of parent commit (NULL for first); UNIQUE enforces linear chain \\
author\_id   & INT          & \fk\ $\to$ USER, NOT NULL     & User who made this commit \\
message      & TEXT         & NOT NULL                      & Commit description message \\
created\_at  & DATETIME     & DEFAULT CURRENT\_TIMESTAMP   & Commit creation timestamp \\
\bottomrule
\end{tabular}
\end{center}

\subsubsection*{Merkle Hash Construction}

The \texttt{hash} is not a random UUID --- it is a deterministic cryptographic
fingerprint computed as:

\[
\texttt{hash} = \text{SHA-256}\bigl(\,\text{serialise}(\textit{schema},\,
\textit{data})\;\|\;\texttt{parent\_hash}\,\bigr)
\]

This means any modification to a historical snapshot immediately invalidates
all subsequent commit hashes, making unauthorised history rewriting detectable.

\subsubsection*{Design Notes}
\begin{itemize}[leftmargin=1.5em]
    \item \texttt{parent\_hash} is \texttt{NULL} only for the very first commit
    in a repository (the root commit), analogous to Git's initial commit.
    \item Because there is no branching, \texttt{parent\_hash} is unique across
    the table --- no two commits share the same parent. This is now enforced
    at the DDL level via \texttt{UNIQUE KEY uq\_parent\_hash (parent\_hash)},
    guaranteeing the linear chain constraint at the database layer.
    \item The self-referencing foreign key on \texttt{parent\_hash} enforces
    referential integrity of the chain at the database level.
\end{itemize}

% ── SNAPSHOTS ─────────────────────────────────
\subsection{SNAPSHOTS}

The \texttt{SNAPSHOTS} entity is the largest and most storage-intensive table in
\texttt{gitdb\_meta}. It stores the \textit{complete state} of every tracked
table at the time of each commit --- both the schema (DDL) and the data (rows)
--- as JSON blobs. This is what enables GitDB's state-based approach: rather than
storing migration instructions, it stores the actual state.

\begin{center}
\begin{tabular}{>{\ttfamily}l >{\ttfamily}l l l}
\toprule
\textbf{Attribute} & \textbf{Type} & \textbf{Constraint} & \textbf{Description} \\
\midrule
id           & BIGINT       & \pk, AUTO\_INCREMENT          & Surrogate primary key \\
commit\_hash & CHAR(64)     & \fk\ $\to$ COMMITS, NOT NULL  & Commit this snapshot belongs to \\
table\_name  & VARCHAR(255) & NOT NULL                      & Name of the MySQL table snapshotted \\
ddl\_json    & LONGTEXT     & NOT NULL                      & Serialised DDL (schema) as JSON \\
rows\_json   & LONGTEXT     & NOT NULL                      & Serialised table rows as JSON array \\
row\_count   & INT          & NOT NULL                      & Number of rows at snapshot time \\
captured\_at & DATETIME     & DEFAULT CURRENT\_TIMESTAMP   & Timestamp of snapshot capture \\
\bottomrule
\end{tabular}
\end{center}

\subsubsection*{JSON Storage Format}

\textbf{\texttt{ddl\_json}} stores column metadata extracted from
\texttt{information\_schema.COLUMNS}:
\begin{lstlisting}[style=sqlstyle, caption={Example ddl\_json structure}]
{
  "table_name": "orders",
  "columns": [
    { "name": "order_id",   "type": "INT",          "nullable": false,
      "key": "PRI", "default": null,  "extra": "auto_increment" },
    { "name": "customer_id","type": "INT",          "nullable": false,
      "key": "MUL", "default": null,  "extra": "" },
    { "name": "total",      "type": "DECIMAL(10,2)","nullable": false,
      "key": "",    "default": "0.00","extra": "" }
  ]
}
\end{lstlisting}

\textbf{\texttt{rows\_json}} stores a JSON array of all rows, ordered by primary
key for deterministic hashing:
\begin{lstlisting}[style=sqlstyle, caption={Example rows\_json structure}]
[
  { "order_id": 1, "customer_id": 42, "total": "149.99" },
  { "order_id": 2, "customer_id": 17, "total":  "89.50" }
]
\end{lstlisting}

\subsubsection*{Design Notes}
\begin{itemize}[leftmargin=1.5em]
    \item A \texttt{UNIQUE} constraint on \texttt{(commit\_hash, table\_name)}
    ensures no duplicate snapshot exists for the same table in the same commit.
    \item \texttt{LONGTEXT} is used instead of \texttt{JSON} type to maximise
    compatibility and avoid MySQL JSON indexing overhead for blobs that are only
    read and written atomically.
    \item A hard limit of \textbf{10,000 rows per table} is enforced at commit
    time to keep snapshot sizes and in-memory diffing within practical bounds.
    \item Tables without a primary key cannot be diffed and are flagged at
    commit time with a diagnostic message.
\end{itemize}

% ─────────────────────────────────────────────
\section{Relationships}

\begin{center}
\begin{tabular}{l l l p{6.5cm}}
\toprule
\textbf{From} & \textbf{To} & \textbf{Name} & \textbf{Cardinality \& Description} \\
\midrule
USER       & REPOSITORY & \textit{owns}      & 1:N --- One user owns many repositories; each repository has exactly one owner \\[4pt]
USER       & COMMITS    & \textit{authors}   & 1:N --- One user authors many commits; each commit has exactly one author \\[4pt]
REPOSITORY & COMMITS    & \textit{contains}  & 1:N --- One repository contains many commits; each commit belongs to exactly one repository \\[4pt]
COMMITS    & COMMITS    & \textit{child of}  & 1:0..1 (self) --- Each commit optionally has one parent and optionally has one child; no branching means the chain is strictly linear \\[4pt]
COMMITS    & SNAPSHOTS  & \textit{has}       & 1:N --- One commit has many snapshots (one per tracked table); each snapshot belongs to exactly one commit \\[4pt]
REPOSITORY & COMMITS    & \textit{points to} & 1:1 via \texttt{current\_hash} --- Each repository's HEAD points to exactly one commit (NULL before the first commit) \\
\bottomrule
\end{tabular}
\end{center}

% ─────────────────────────────────────────────
\section{Complete DDL}

The following SQL creates the complete \texttt{gitdb\_meta} schema including
all constraints.

\begin{lstlisting}[style=sqlstyle, caption={gitdb\_meta complete DDL}]
CREATE DATABASE IF NOT EXISTS gitdb_meta;
USE gitdb_meta;

-- Users
CREATE TABLE user (
    user_id       INT           NOT NULL AUTO_INCREMENT,
    username      VARCHAR(50)   NOT NULL,
    email         VARCHAR(100)  NOT NULL,
    password_hash VARCHAR(255)  NOT NULL,
    full_name     VARCHAR(100)  NOT NULL,
    is_active     BOOLEAN       NOT NULL DEFAULT TRUE,
    created_at    DATETIME      NOT NULL DEFAULT CURRENT_TIMESTAMP,
    PRIMARY KEY (user_id),
    UNIQUE KEY uq_username (username),
    UNIQUE KEY uq_email    (email)
);

-- Repositories
CREATE TABLE repository (
    repo_id          INT           NOT NULL AUTO_INCREMENT,
    user_id          INT           NOT NULL,
    repo_name        VARCHAR(100)  NOT NULL,
    target_db_name   VARCHAR(100)  NOT NULL,
    db_host          VARCHAR(100)  NOT NULL,
    db_port          INT           NOT NULL DEFAULT 3306,
    db_user          VARCHAR(50)   NOT NULL,
    db_password_key  VARCHAR(100)  NULL,     -- keyring key, not the password
    current_hash     CHAR(64)      NULL,
    created_at       DATETIME      NOT NULL DEFAULT CURRENT_TIMESTAMP,
    PRIMARY KEY (repo_id),
    CONSTRAINT fk_repo_user
        FOREIGN KEY (user_id) REFERENCES user (user_id)
);

-- Commits (Merkle chain)
CREATE TABLE commits (
    hash         CHAR(64)      NOT NULL,
    repo_id      INT           NOT NULL,
    parent_hash  CHAR(64)      NULL,
    author_id    INT           NOT NULL,
    message      TEXT          NOT NULL,
    created_at   DATETIME      NOT NULL DEFAULT CURRENT_TIMESTAMP,
    PRIMARY KEY (hash),
    UNIQUE KEY uq_parent_hash (parent_hash),   -- enforces linear chain at DB level
    CONSTRAINT fk_commit_repo
        FOREIGN KEY (repo_id)     REFERENCES repository (repo_id),
    CONSTRAINT fk_commit_parent
        FOREIGN KEY (parent_hash) REFERENCES commits    (hash),
    CONSTRAINT fk_commit_author
        FOREIGN KEY (author_id)   REFERENCES user       (user_id)
);

-- Add HEAD FK on repository after commits table exists
ALTER TABLE repository
    ADD CONSTRAINT fk_repo_head
        FOREIGN KEY (current_hash) REFERENCES commits (hash);

-- Snapshots
CREATE TABLE snapshots (
    id           BIGINT        NOT NULL AUTO_INCREMENT,
    commit_hash  CHAR(64)      NOT NULL,
    table_name   VARCHAR(255)  NOT NULL,
    ddl_json     LONGTEXT      NOT NULL,
    rows_json    LONGTEXT      NOT NULL,
    row_count    INT           NOT NULL,
    captured_at  DATETIME      NOT NULL DEFAULT CURRENT_TIMESTAMP,
    PRIMARY KEY (id),
    UNIQUE KEY uq_commit_table (commit_hash, table_name),
    CONSTRAINT fk_snap_commit
        FOREIGN KEY (commit_hash) REFERENCES commits (hash)
);
\end{lstlisting}

% ─────────────────────────────────────────────
\section{Storage and Scale Considerations}

\begin{center}
\begin{tabular}{l l p{8cm}}
\toprule
\textbf{Constraint} & \textbf{Value} & \textbf{Rationale} \\
\midrule
Max rows per table  & 10,000   & Keeps \texttt{rows\_json} blobs under $\sim$5\,MB; ensures in-memory diff completes in practical time \\
Snapshot storage    & LONGTEXT & Supports up to 4\,GB per cell; sufficient for all mini-project workloads \\
Hash length         & CHAR(64) & Fixed-length SHA-256 hex digest; no variable-length overhead \\
Snapshot uniqueness & UNIQUE (commit\_hash, table\_name) & Prevents duplicate snapshots for the same table in the same commit \\
\bottomrule
\end{tabular}
\end{center}

\medskip
\noindent For production use beyond mini-project scope, \texttt{rows\_json}
could be compressed (e.g.\ \texttt{ZLIB} via application layer) or replaced
with a dedicated object store, with only a content-addressed reference stored
in \texttt{gitdb\_meta}.

\end{document}
\end{antml:parameter>